% D.77 candidate. UNREGISTERED, NOTES ONLY.
% Written during checkpoint11 source freeze. No inclusion or promotion is authorized.
\subsection{The entire contribution at five is peripheral}

\begin{statement}{D.77}{Open}{Candidate: exact local and global invariant factors}
Let \(s=(a,b)\) be any completed abelian Kummer pair for the tripod over
\(\mathbb Q\), with arbitrary ramification, and use the actual Alexander
module \(A_s\) of D.35. Retain \(L_5,\mathcal T_s,\mathcal S_s,\mathcal K_s\)
from D.63. Set
\[
 (c_1,c_2,c_3)=(a,b,a-b).
\]
For \(c\ne0\), let \(r(c)\) and \(N(c)\) be the global and local-at-five
divisibility exponents of D.75. For \(c=0\), set \(r(c)=N(c)=\infty\).
Write
\[
 Z(d)=
 \begin{cases}
 \mathbb Z_5^d,&d<\infty,\\
 \displaystyle\prod_{j\ge1}\mathbb Z_5,&d=\infty .
 \end{cases}
 \qquad Z(0)=0.                                      \tag{D.77.1}
\]
The candidate gives isomorphisms of compact modules
\[
 L_5\simeq\prod_{i=1}^3 Z(N(c_i)),\qquad
 \mathcal T_s\simeq\prod_{i=1}^3 Z(r(c_i)).             \tag{D.77.2}
\]
They may be chosen so that \(L_5\to\mathcal T_s\) is, on each nonzero
axis, the projection to its first \(r(c_i)\) coordinates. On a zero
axis it is the identity between the two countable products.

Define \(\delta(c)\) as in D.76, with \(\delta(0)=0\), and put
\(d(s)=\delta(a)+\delta(b)+\delta(a-b)\). Then the candidate identifies
the \emph{whole} kernel, not just a split factor, as
\[
 \ker(L_5\longrightarrow\mathcal T_s)\simeq Z(d(s)),
 \qquad
 \mathcal K_s\simeq\mathcal S_s\times Z(d(s)).          \tag{D.77.3}
\]
The latter splitting is continuous and noncanonical.
Here \(\mathcal S_s\) is exactly the inverse limit of the finite
global localization kernels in (D.63.5). It is not asserted to vanish
or to be an elliptic-curve Tate--Shafarevich group.
This candidate makes no assertion that a section exists above \(s\),
and does not compute a first higher-derived obstruction.
\end{statement}

\begin{proof}[Candidate proof, pending independent audits]
\paragraph{A primitive local Soul\'e character without a special-value input.}
Fix the embedding \(G_{\mathbb Q_5}\subset G_{\mathbb Q}\) and put
\[
 K=\mathbb Q_5(\zeta),\quad \zeta=\zeta_5,\quad
 \pi=\zeta-1,\quad
 \phi=-\zeta^2-\zeta^3=(1+\sqrt5)/2.
\]
The polynomial
\[
 \Phi_5(1+\pi)=\pi^4+5\pi^3+10\pi^2+10\pi+5
                                                               \tag{D.77.4}
\]
is Eisenstein. Normalize \(v_\pi(\pi)=1\), so \(v_\pi(5)=4\),
and let \(U_j=1+\pi^j\mathcal O_K\). The identity
\[
 \phi^4-1=3\phi+1=(5+3\sqrt5)/2
                                                               \tag{D.77.5}
\]
has valuation exactly two. Every fifth power in \(U_1\) lies in
\(U_5\): in \((1+x)^5-1\), the first four terms have valuations
at least \(4+jv_\pi(x)\ge5\), and the last has valuation at
least five. Also \(v_\pi(\zeta^e-1)=1\) for \(5\nmid e\).

These facts show that \(\phi,\zeta\) are independent in
\(K^\times/K^{\times5}\). Indeed, suppose
\(\phi^m\zeta^e=x^5\), where \(0\le m,e<5\).
The element \(x\) is a unit. Its fourth power is in \(U_1\),
since the residue field is \(\mathbb F_5\). Hence
\(\phi^{4m}\zeta^{4e}\in U_5\). If \(e\ne0\), its difference
from one has valuation one: the \(\phi^{4m}\) factor is in
\(U_2\), whereas \(\zeta^{4e}\) is not. If \(e=0\) and
\(m\ne0\), its difference from one has valuation two, by
(D.77.5) and \(5\nmid m\). Both are contradictions.

The raw Soul\'e normalization is (D.39.1), equivalently
\cite[(1.1)--(1.2)]{DE_NakamuraSakugawaWojtkowiak2017}.
D.35's exact first-level identity is
\[
 u_{3,1}=\prod_{j=1}^4(1-\zeta^j)^{j^2}
        =\phi^3\bigl((1-\zeta)^6\phi^2\bigr)^5.
                                                               \tag{D.77.6}
\]
The cyclotomic extension \(K(\mu_{5^\infty})/K\) has Galois group
\(\mathbb Z_5\), and its unique degree-five subextension is
\(K(\zeta_{25})\), with Kummer line generated by \(\zeta\).
The independence just proved implies
\(K(\phi^{1/5})\not\subset K(\mu_{5^\infty})\).
Consequently, if
\[
 \Delta_5=\ker\{\chi:G_{\mathbb Q_5}\to\mathbb Z_5^\times\},
\]
then \(\kappa_3\bmod5\) is nonzero on \(\Delta_5\).
In particular its actual integral image there contains a unit.
This proves the required local primitivity directly. No global
class-number inference or local \(L\)-value nonvanishing is used.

\paragraph{An actual local element annihilates the interior.}
On \(\Delta_5\), the three maps \(a,b,\kappa_3\) are continuous
additive homomorphisms. Their image
\[
 M=\{(a(\tau),b(\tau),\kappa_3(\tau)):\tau\in\Delta_5\}
       \subset\mathbb Z_5^3
\]
is a compact additive subgroup, and hence a closed
\(\mathbb Z_5\)-submodule. The Kummer and Soul\'e covariance gives,
under conjugation by an element with cyclotomic value \(\omega\),
the operator
\[
 T=\operatorname{diag}(\omega,\omega,\omega^3),
 \qquad \omega^2=-1,\quad\omega\in\mathbb Z_5^\times.
\]
Such an actual element of \(G_{\mathbb Q_5}\) exists, since its
cyclotomic character is onto. It is not required to have order four.
The integral operator
\[
 \frac{T-\omega}{\omega^3-\omega}
             =\operatorname{diag}(0,0,1)                \tag{D.77.7}
\]
preserves \(M\), because its denominator is a five-adic unit.
Apply it to an element whose last coordinate is a unit.
By the definition of the actual image \(M\), there is therefore
an actual \(\delta\in\Delta_5\) such that
\[
 a(\delta)=b(\delta)=0,\qquad \kappa_3(\delta)\in\mathbb Z_5^\times .
                                                               \tag{D.77.8}
\]
No formal assignment of Galois coefficients has occurred.

Put \(t_1=1+u,\ t_2=1+v,\ w=t_1t_2-1\), and use the actual beta
series \(\mathcal B_g\) of D.39. On \(\delta\), its exact boundary
values are
\[
 \mathcal B_\delta(t,1)=\mathcal B_\delta(1,t)
          =\mathcal B_\delta(t,t^{-1})=1.
\]
The first two follow from \(\mathcal B_g=1+uvB_g\), and the last
from the exact reflection formula of D.39 with \(\chi(\delta)=1\).
Thus \(B_\delta\) vanishes under \(v=-u/(1+u)\): the factor
\(uv=-u^2/(1+u)\) is a nonzero divisor in \(\mathbb Z_5[[u]]\).
The kernel of this substitution is precisely \((w)\), by division
by the linear polynomial
\(w=(1+u)(v+u/(1+u))\). Consequently
\[
 \mathcal B_\delta-1=uvw\,E_\delta,\qquad
                E_\delta\in\Lambda^\times.             \tag{D.77.9}
\]
To check the unit assertion, the degree-three term in D.35/D.39 is
\[
 \frac{\kappa_3(\delta)}{48}\,uv(u+v).
 \]
Since \(uvw\) has leading term \(uv(u+v)\), the constant term of
\(E_\delta\) is the unit \(\kappa_3(\delta)/48\).
This coefficient is the raw normalization in (D.77.6), not a
renormalized nonprimitive character.

Let \(D_s=A_s^D\) be the discrete Cartier dual, and let
\(H\) be either \(G_{\mathbb Q}\) or \(G_{\mathbb Q_5}\).
By (D.77.8), the action of \(\delta\) on \(A_s\) is exactly
multiplication by \(\mathcal B_\delta\). Its cyclotomic multiplier
is one, so Cartier twisting adds no scalar in this step.
Every \(H\)-invariant functional therefore annihilates
\((\mathcal B_\delta-1)A_s=uvwA_s\). The ideal \((uvw)\)
is stable, because its three constituent boundary ideals are stable.
With
\[
                         B=A_s/uvwA_s
\]
we have proved
\[
                         D_s^H=(B^D)^H.                \tag{D.77.10}
\]
The same actual local \(\delta\) proves this for the global group.

\paragraph{The finite relation among all three boundaries.}
Let \(M_i=A_s/J_iA_s\), where
\[
                  J_1=(v),\qquad J_2=(u),\qquad J_3=(w).
\]
D.76 proves that the actual modules \(M_i\) are respectively
\(M_a,M_b,M_{a-b}\), with the integral half-cyclotomic gauge on
the last axis. Explicitly, its multiplier before that gauge is
\[
 \chi S_\chi(t)\,t^{\,a(g)-b(g)+(1-\chi(g))/2}.
                                                               \tag{D.77.11}
\]
This follows from D.39's exact reflection identity, so includes
all arithmetic coefficients.

The integral restriction map fits in the exact sequence
\[
 0\longrightarrow B\longrightarrow
        M_1\oplus M_2\oplus M_3
          \longrightarrow T\longrightarrow0,
 \qquad T\simeq\mathbb Z_5^3.                            \tag{D.77.12}
\]
Indeed a triple \((F(u),G(v),H(u))\) is in its image exactly if
\[
 F_0=G_0=H_0,\qquad H_1=F_1-G_1.
\]
After matching the first two axes, any higher diagonal
discrepancy is corrected by \(uvL(u)\), since its diagonal
restriction is \(-u^2L(u)/(1+u)\). The intersection of the
three ideals is \((uvw)\): a multiple \(uvL\) vanishing on the
third axis has \(L\) vanishing there, hence is a multiple of \(uvw\).
This also proves exactness at the first term.

Here the action on the finite free cokernel is explicit. Use its
coordinates
\[
 q_1=G_0-F_0,\qquad q_2=H_0-F_0,\qquad q_3=H_1-F_1+G_1.
\]
For \(\alpha=\chi(g)\), before any individual translation gauges,
the matrix is
\[
 \begin{pmatrix}
 \alpha^2&0&0\\
 0&\alpha^2&0\\
 \alpha^2(b(g)+(\alpha-1)/2)&
 \alpha^2(a(g)-b(g))&\alpha^3
 \end{pmatrix}.                                        \tag{D.77.13}
\]
This comes by expanding the actual multipliers to first degree.
It also gives a filtration of \(T\) with characters
\(\chi^2,\chi^2,\chi^3\), without requiring that the extension
between these characters split.

Let \(D_i=M_i^D\). The compact sequence (D.77.12) splits as
\(\mathbb Z_5\)-modules by choosing lifts of the three basis vectors
of its finite free quotient; these lifts define a continuous map.
In particular dualizing gives the exact sequence
\[
 0\longrightarrow T^D\longrightarrow
       D_1\oplus D_2\oplus D_3
          \longrightarrow B^D\longrightarrow0.         \tag{D.77.14}
\]
The Cartier dual \(T^D\) has successive characters
\(\chi^{-1},\chi^{-1},\chi^{-2}\).
Choose an actual \(k\in H\) with \(\chi(k)=2\). Then \(k-1\) is
an automorphism of \(T^D\). For example the determinant of its
corresponding integral lattice matrix is
\[
             (2^{-1}-1)^2(2^{-2}-1)=-3/16,
                                                               \tag{D.77.15}
\]
a five-adic unit. This conclusion is unaffected by the off-diagonal
entries in (D.77.13).

\paragraph{Every invariant boundary sum has an invariant lift.}
Given \(x\in(B^D)^H\), choose any lift
\(z\in D_1\oplus D_2\oplus D_3\). Then \((k-1)z\in T^D\).
Subtracting its unique inverse image under \(k-1\) produces
a lift \(z'\) fixed by \(k\).

Use the measure model of D.75 on each axis. Its elements are
continuous functions on \(\mathbb Z_5\) with values in
\(\mathbb Q_5/\mathbb Z_5\), modulo constants. An integral
translation makes the affine action of \(k\) pure multiplication
by two, since \(\chi(k)-1=1\). This includes the integral gauge
in (D.77.11). Represent each component of \(z'\) by \(f_i\),
with \(f_i(0)=0\). Being fixed by \(k\) says
\[
                              f_i(2z)=f_i(z).           \tag{D.77.16}
\]
The possible constant defect is zero by evaluation at zero.

For \(\tau\in\ker(\chi|_H)\), the defect \(\tau z'-z'\)
belongs to \(T^D\), by invariance of \(x\). A component of \(T^D\)
involves only the constant and linear coefficients of its axis
power series. Under the augmentation-ideal measure model,
\[
 \frac{t^z-1}{t-1}=z+\binom z2(t-1)+\cdots,
\]
so that component is represented by a polynomial of degree at
most two in \(z\), modulo constants. Translation gauges preserve
this degree bound. Hence, with \(h=c_i(\tau)\), changing the sign
of \(h\) if required by the dual-action convention,
\[
 f_i(z+h)-f_i(z)=P_h(z)=A_hz^2+B_hz+C_h.                 \tag{D.77.17}
\]
The coefficients are unique in \(\mathbb Q_5/\mathbb Z_5\):
evaluation at \(0,1,2\) has determinant two. In particular
coefficient comparison is legitimate for this torsion target.

Equation (D.77.16) and telescoping give
\[
 P_{2h}(2z)=P_h(z),\qquad P_{2h}(z)=P_h(z+h)+P_h(z).
\]
The quadratic terms imply \(A_{2h}=2A_h\) and
\(4A_{2h}=A_h\), hence \(7A_h=0\) and \(A_h=0\).
The linear terms then imply \(B_{2h}=2B_h\) and
\(2B_{2h}=B_h\), hence \(3B_h=0\) and \(B_h=0\).
Finally \(C_{2h}=2C_h=C_h\), so \(C_h=0\).
Thus every component of \(z'\) is fixed by \(\ker(\chi|_H)\).

The element two topologically generates \(\mathbb Z_5^\times\):
its order modulo five is four, and \(v_5(2^4-1)=1\), so its
fourth power generates \(1+5\mathbb Z_5\).
Consequently the closed subgroup generated by \(k\) and
\(\ker(\chi|_H)\) is all of \(H\). The components of \(z'\)
are therefore \(H\)-invariant. Conversely the invariant kernel
in (D.77.14) is zero, since \(k-1\) is invertible on \(T^D\).
Together with (D.77.10), this proves the natural isomorphism
\[
 D_1^H\oplus D_2^H\oplus D_3^H
                  \xrightarrow{\ \sim\ }D_s^H
       \quad(H=G_{\mathbb Q},G_{\mathbb Q_5}).           \tag{D.77.18}
\]
The map is the actual sum of the three inclusions, so these two
isomorphisms commute with the inclusion of global in local
invariants. No vanishing of \(H^1(H,T^D)\) has been assumed.

\paragraph{All local and global invariant factors and the remaining term.}
D.75 computes each group in (D.77.18). For a depth \(d<\infty\),
it is \((\mathbb Q_5/\mathbb Z_5)^d\), with basis the functions
\(\mathbf1_{5^j\mathbb Z_5}\), \(1\le j\le d\), modulo constants.
For the zero local class, \(d=\infty\), it is the direct sum of
these groups over all \(j\ge1\): continuity forces finite support
in the ball basis. The compatible global and local gauges in
D.75 identify the first \(r(c_i)\) vectors under restriction
when \(c_i\ne0\). If \(c_i=0\), one global gauge makes both
actions pure unit actions and the inclusion is the identity.

Local duality and the finite-coefficient inverse limits of D.63
identify
\[
 L_5=(D_s^{G_{\mathbb Q_5}})^*,\qquad
 \mathcal T_s=(D_s^{G_{\mathbb Q}})^*.
                                                               \tag{D.77.19}
\]
The star here is ordinary Pontryagin duality of discrete torsion
groups. In particular it is not an additional Cartier twist.
The finite local cohomology groups are finite, so their inverse
systems are Mittag--Leffler; invariants commute with the increasing
union of finite duals. These are exactly the identifications
proved in D.63 from \cite[I.2.1, I.2.3, I.4.10]{DE_MilneADT}.
Thus this step does not assume that one finite-level class
survives to an infinite compatible system.

Dualizing (D.77.18), with the explicit ball bases, proves
(D.77.2) and its assertion about the actual map. Its entire
kernel is \(Z(d(s))\). D.63 now supplies the exact sequence
\[
 0\longrightarrow\mathcal S_s\longrightarrow\mathcal K_s
       \longrightarrow Z(d(s))\longrightarrow0.
                                                               \tag{D.77.20}
\]
The discrete dual of the last term is a direct sum of copies of
\(\mathbb Q_5/\mathbb Z_5\), and is divisible. It is therefore
injective as a \(\mathbb Z_5\)-module: the extension criterion
for ideals \(5^n\mathbb Z_5\) is exactly divisibility by \(5^n\).
Its inclusion in the discrete dual of \(\mathcal K_s\) has a
retraction. Pontryagin duality turns that retraction into a
continuous section of (D.77.20), proving (D.77.3).
Only the compact groups in D.63 are split here. No splitting
of the original nonseparated global cochain quotient is asserted.

The executed exact certificate
\texttt{notes/de\_local\_boundary\_certificate.py} verifies the 24
nontrivial Kummer combinations in the Eisenstein ring, the raw
Soul\'e unit factorization, the image projector, the actual jet
matrix and determinant \(-3/16\), and the order twenty of two
modulo twenty-five. Its JSON and output are retained. These finite
identities support the integral arguments above; they do not
replace the infinite invariant and duality proofs.
\end{proof}
